% Einheit 1 von 4 – Das Skalarprodukt als Objekt (bank/skalarprodukt-und-winkel/e1.jsonl, zusammenbau v0.1)
%% TODO Zeitmarke Einheit 1: Marken-Zeile nicht lesbar
%% TODO Zweigzeile Teil 1: was hier gelernt wird (Satz in Schülersprache aus dem Einheitstitel) – Einheit 1
\einheitenkopf[e1]{Einheit 1 von 4 · Das Skalarprodukt als Objekt}
\zweigzeile{keine Prüfungsaufgabe}
\setcounter{aufgabe}{8}

%% TODO Titel: Ich-kann-Satz fehlt in der Bank (Platzhalter: Kettenname) – Nr. 9
%% TODO Anweisung: ein Satz über der ersten Teilaufgabe fehlt in der Bank – Nr. 9
\begin{aufgabe}{Das Skalarprodukt als Objekt}
%% TODO \rechenplatz{n}: Schrittzahl des Grundfalls fehlt in der Bank (nur bei fester Schreibform, 2.3 b)
\begin{teile}
\teil $\vec a$, $\vec b$ und $\vec c$ sind Vektoren mit je drei Koordinaten. Was ergibt $\vec a \circ \vec b$? Kreuze an. \\ \kreuz{Vektor} \\ \kreuz{Zahl} \\ \kreuz{nicht definiert}
\teil $\vec a = (2 \mid 1 \mid 3)$, $\vec b = (1 \mid 4 \mid -1)$. Berechne $\vec a \circ \vec b$ und gib an, ob der Winkel zwischen $\vec a$ und $\vec b$ spitz, recht oder stumpf ist. $\vec a \circ \vec b =$ \leerfeld \quad Winkel: \leerfeld
\teil $\vec a$, $\vec b$ und $\vec c$ sind Vektoren mit je drei Koordinaten. Prüfe die Rechenarten von innen nach außen: Was ergibt $(\vec a + \vec b) \circ \vec c$? \\ \kreuz{Vektor} \\ \kreuz{Zahl} \\ \kreuz{nicht definiert}
\teil $\vec u = (1 \mid 2 \mid -1)$ und $\vec v = (r \mid 3 \mid 4)$ mit $r \in \mathbb{R}$. Für welche Werte von r ist der Winkel zwischen $\vec u$ und $\vec v$ mindestens $90^\circ$ groß? \leerfeld
\teil Betrachtet werden alle Vektoren $(x \mid y \mid z)$ mit nichtnegativen Koordinaten und $x + y + z = 10$. Beurteile die Aussage: „Es gibt zwei verschiedene solche Vektoren, die den Winkel $0^\circ$ einschließen.“
\teil Ein Quader hat eine quadratische Grundfläche mit der Seitenlänge s und die Höhe h. An einer Ecke der Grundfläche sind $\vec u$ und $\vec v$ die Kantenvektoren der Grundfläche und $\vec w$ der Kantenvektor der Höhe. Begründe ohne Koordinaten, dass $\vec u \circ (\vec u + \vec v + 2\vec w)$ nur von s abhängt. \hfill (Abitur 2021 GK)
\end{teile}
\end{aufgabe}

%% TODO Titel: Ich-kann-Satz fehlt in der Bank (Platzhalter: Kettenname) – Nr. 10
%% TODO Anweisung: ein Satz über der ersten Teilaufgabe fehlt in der Bank – Nr. 10
\begin{aufgabe}{Das Skalarprodukt als Objekt – Fehler finden, Begründen}
\begin{teile}
\teil $\vec a$, $\vec b$ und $\vec c$ sind Vektoren mit je drei Koordinaten. Jan soll angeben, was $\vec a + (\vec b \circ \vec c)$ ergibt, und kreuzt „Vektor“ an, weil $\vec a$ ein Vektor ist. Finde den Fehler.
\teil Begründe, warum $\vec a \circ \vec b$ eine Zahl und kein Vektor ist.
\end{teile}
\end{aufgabe}
